\subsection{Comparison of the inverse data}
\label{sec:inverse-data-comparison}

Classical inverse constructions already recover a web from suitable
Reye data.  In \cite[Proposition 3.6]{MMV}, a Fano polarization
\(H\) determines its Reye bundle, when that bundle exists.
Theorem 3.7 of that work realizes the polarized Enriques surface in
a Grassmannian.  In characteristic different from two, Remark 3.8
uses the canonical K3 cover, its involution, and two exchanged line
bundles \(L^+,L^-\).  A four-dimensional multiplication kernel
defines symmetric \((1,1)\)-equations and hence a web.  Thus
reconstruction from an Enriques polarization with its Reye structure
is not asserted here as new.  Nor is every nodal Enriques surface a
regular classical Reye web: the distinction is explicit in
\cite[Remark 1.2 and Section 6]{MMV}.

These inputs differ from a single quartic polarization.  The first-layer
invariant in Theorem~\ref{thm:main-reconstruction} recovers
\((Y_R,\OO_{Y_R}(1))\), without assuming the Reye involution or
the second exchanged polarization as marked data.  Theorem
\ref{thm:torelli-monodromy} therefore retains the function-field
degree of this particular unmarked map; the classical reconstruction
with additional data does not compute that degree.  Theorem
\ref{thm:generic-intrinsic-web} starts from neither a marked Reye
surface nor a specified web.  It starts from the abstract nonreduced
Fitting scheme, recovers an oriented deepest normal cone, and extracts
two specified coefficient lines through
Lemma~\ref{lem:universal-schur-readout}.  Lemma
\ref{lem:common-g-functoriality} proves that an abstract isomorphism
acts on both lines through the same projective transformation.
The exterior-support theorem then makes recombination unique on the
entire smooth Jacobian locus, without an additional genericity
restriction.  This input and this extraction mechanism are the inverse
assertions proved here; classical web-to-Reye and web-to-Jacobian
constructions are used as comparisons, not rebranded as new results.

\subsection{Role of the relative primary theorem}
\label{sec:relative-tool-positioning}

The relative algebra supplies a simultaneous model for a bounded
principal-colon tower.  Generic freeness, flattening, associated-point
technology, and effective descent are standard ingredients
\cite{StacksGenericFree,StacksFlattening,StacksRelativeAssassin,StacksDescent}.
Our use of them is a finite-diagram construction, not a new principle
of generic freeness or primary existence.  The outputs needed here
are simultaneous formation of the colons and their quotients after
base change, exhaustion of the geometric associated points, closed
support packets on the original strata, and relative torsion modules
realizing the packet lengths.  The diagonal injections, their
cokernels, and the regular-element witnesses are kept in the same
diagram so that these outputs hold together.  Individual flattening
of prime supports would not imply this combination, and no theorem
about canonical labelled primary decompositions is asserted.
The bounded-principal-colon theorem is therefore a relative-algebra
tool for the geometric results, rather than the basis for an independent
foundational-priority claim.  The determinant identities, calculated
boundary schemes, and intrinsic inverse theorem carry the geometric
content.

\subsection{Extent of the primary tables}
\label{sec:primary-table-scope}

The primary decompositions computed in full are the stated generic
corank-two coefficient strata, their specified collisions, and the
quotient-induced benchmark families.  The first two minimal-support
laws and the sharp index-five nilpotence theorem apply globally on
\(G_4^\circ\).  All their proofs and coefficient tables are retained.
The seven-component two-parameter incidence calculation describes the
specified pullback of an incidence resolution; it does not identify
all embedded primary components of \(W_3\) and \(W_4\) at coranks
three and four.  The inverse theorem does not require such an atlas:
it reads the degree-twelve maximal-minor space itself from the
degree-sixteen ideal, before choosing a primary decomposition.
Accordingly its hypothesis and conclusion are unchanged by this
separation of the two questions.
